% !TeX program = pdflatex
% =============================================================================
% Aufgabenblatt 3: Schaltungen
% UZH Praktikum I -- Sara Romer Urban, Kantonsschule im Lee, HS2026
% Klasse: 1f (23.09.2026), aus PLAN_Dokumentengenerierung.md.
%
% ENTWURF. Uebungsaufgaben zu arbeitsblatt4-schaltungen.tex im selben
% Ordner (Serie-/Parallelschaltung, Ersatzwiderstand, Mischschaltungen,
% Kirchhoffsche Regeln, UND-/ODER-Schaltung) -- Aufteilung nach Loris'
% eigenem Kriterium (Arbeitsblatt = Plenum/Praktisches, Aufgabenblatt =
% Uebungsaufgaben), analog zu u-i-r/aufgabenblatt-u-i-r.tex. Deshalb
% bewusst OHNE einleitende Erklaerbox, beginnt direkt mit der ersten
% Aufgabe.
%
% -----------------------------------------------------------------------------
% Korrektur (24.09.2026, auf Wunsch): Nummerierung (Aufgabenblatt-Zaehlung
% nach institutionen/UZH-Praktikum-I/CLAUDE.md, Abschnitt "Nummerierung
% und Dateibenennung") jetzt festgelegt: Aufgabenblatt 3, passend zum
% bereits vergebenen Arbeitsblatt 4 (arbeitsblatt4-schaltungen.tex im
% selben Ordner) -- die Nummernpaarung Arbeitsblatt = Aufgabenblatt + 1
% ist bei jedem anderen Themenpaar in diesem Repo so belegt (z. B.
% Kondensatoren: Arbeitsblatt 4 / Aufgabenblatt 3, siehe CLAUDE.md).
% Vorwissen/ im elektrizitaet-Ordner enthaelt keine eindeutig nummerierten
% Dokumente Sara Romers zu diesem Thema, die Nummer folgt daher aus dieser
% Paarungskonvention statt direkt aus Vorwissen/ abgeleitet zu sein.
% Datei entsprechend von aufgabenblatt-schaltungen.tex auf
% aufgabenblatt3-schaltungen.tex umbenannt (Titel/PDF-Metadaten ebenso
% angepasst).
%
% Aufgaben 1-4 aus institutionen/UZH-Praktikum-I/arbeitsblaetter/
% elektrizitaet/schaltungen/Aufgaben/ (Loris' eigene Sammlung, Quelle
% LEIFIphysik.de "Komplexere Schaltkreise") uebernommen und angepasst:
% eigene Formulierung/eigene circuitikz-Schaltplaene statt der
% Original-Screenshots, Teilaufgaben zu Leistung/Energie (P=U*I) entfernt,
% da im Arbeitsblatt nicht behandelt -- nur Widerstand/Spannung/Strom
% (Ohmsches Gesetz, Ersatzwiderstand) bleibt im Rahmen des Gelernten.
% - Aufgabe 1 (qualitativ, Serie/Parallel-Vergleich): aus Aufgaben/
%   Aufgabe_1 "Zweite Lampe im Stromkreis", unveraendert uebernommen
%   (bereits rein qualitativ, kein Leistungsbezug).
% - Aufgabe 2 ((R1+R2) parallel zu R3, U0=220V): aus Aufgaben/Aufgabe_2
%   "Widerstandskombination", nur Teilaufgaben a) Rges und b) Stroeme
%   uebernommen (Teilaufgaben c/d zur Leistung weggelassen).
% - Aufgabe 3 ((R1 parallel R2) in Serie mit R3, U0=12V): aus Aufgaben/
%   Aufgabe_0 "Minimaler und maximaler Widerstand", vollstaendig
%   uebernommen (a: qualitative Abschaetzung, b: Strom, c: Teilspannungen
%   -- reine Ersatzwiderstand-/Ohmsches-Gesetz-Anwendung, kein
%   Leistungsbezug).
% - Aufgabe 4 ((R1 parallel R2) in Serie mit (R3 parallel R4), U=12V):
%   aus Aufgaben/Aufgabe_3 "Widerstandsaenderung", vollstaendig
%   uebernommen (a: Rges und Strom bei A, b: qualitative Folgeueberlegung
%   bei groesserem R1).
% (Aufgabe 5, urspruenglich UND-/ODER-Schaltung, siehe Korrektur weiter
% unten -- inzwischen wieder entfernt.)
%
% -----------------------------------------------------------------------------
% Review-Korrektur (22.09.2026, auf Wunsch -- gemeinsam mit
% arbeitsblatt4-schaltungen.tex per physics-accuracy-reviewer geprueft):
% Aufgabe 3c verwendete den in 3b auf 0.102 A gerundeten Strom weiter,
% was zu einem kleinen Rundungsfehler bei $U_3$ fuehrte (8.16 V statt der
% praeziseren 8.17 V). Zwischenwert in 3b auf 0.1021 A praezisiert, 3c
% entsprechend auf 8.17 V/3.83 V angepasst. Ausserdem stellte sich beim
% Review heraus, dass Aufgabe 3a implizit die allgemeine Regel
% "Ersatzwiderstand einer Parallelschaltung ist immer kleiner als der
% kleinste Einzelwiderstand" voraussetzt -- diese stand damals in einer
% Merksatz-Box in arbeitsblatt4-schaltungen.tex sichtbar; diese Box ist
% seither wieder entfernt worden (siehe Korrektur unten), die Regel steckt
% jetzt in der Konklusion-Antwort, die die SuS im Arbeitsblatt selbst zu
% dieser Schaltung herleiten.
%
% -----------------------------------------------------------------------------
% Korrektur (22.09.2026, auf Wunsch): (1) Aufgabe 5 (UND-/ODER-Schaltung)
% komplett entfernt: Sie setzte die Begriffe "UND-Schaltung"/
% "ODER-Schaltung" voraus, die im Arbeitsblatt aber nur in der optionalen
% "Zusatzaufgabe (fuer schnelle Gruppen)" vorkommen. Nicht jede Gruppe
% bearbeitet diese, die Begriffe sind also nicht zuverlaessig eingefuehrt.
% Das Aufgabenblatt hat dadurch nur noch 4 statt 5 Aufgaben (\setcounter
% -Werte entsprechend unveraendert, da Aufgabe 5 ohnehin die letzte war).
% (2) Die selbst gezeichneten circuitikz-Schaltplaene in Aufgabe 2, 3 und 4
% durch die tatsaechlichen Referenzbilder aus den jeweiligen Aufgaben/
% -Unterordnern ersetzt (von SVG nach PDF konvertiert via Inkscape, siehe
% gespeicherte Notiz zum SVG-Workflow): Aufgabe_2/
% widerstandskombination-schaltplan.pdf, Aufgabe_0/
% minimaler-maximaler-widerstand-schaltplan.pdf (fuer Aufgabe 3) und
% Aufgabe_3/widerstandsaenderung-schaltplan.pdf (fuer Aufgabe 4):
% inhaltlich identisch zu den bisherigen eigenen Schaltplaenen, aber die
% Originalvorlage statt einer Nachzeichnung. Aufgabe 1 hat kein SVG im
% zugehoerigen Aufgaben/Aufgabe_1/-Ordner (nur die LEIFI-PDF), ihr
% einfacher Ein-Gluehbirne-Schaltplan bleibt deshalb unveraendert als
% eigene circuitikz-Zeichnung.
% =============================================================================
\documentclass[addpoints,11pt,a4paper]{exam}

\usepackage[T1]{fontenc}
\usepackage[utf8]{inputenc}
\usepackage[ngerman]{babel}
\usepackage{lmodern}
\usepackage[a4paper,margin=2.2cm,headheight=14pt]{geometry}
\usepackage{amsmath}
\usepackage{siunitx}
% Hausstil: zusammengesetzte Einheiten immer als echter Bruch (\frac), nicht
% als Inline-Schraegstrich oder \tfrac -- gilt global fuer jedes \si/\SI.
\sisetup{per-mode=fraction,fraction-command=\frac}
\usepackage{array,booktabs,tabularx,multirow}
\usepackage{enumitem}
\usepackage{xcolor}
\usepackage{colortbl}
\usepackage{tikz}
\usepackage[european]{circuitikz}
\usepackage[most]{tcolorbox}
\usepackage{lastpage}
\usepackage{graphicx}
\usepackage{hyperref}

\usetikzlibrary{arrows.meta,calc,angles,quotes}

\usepackage{setspace}
\usepackage{needspace}
\setlength{\parindent}{0pt}
\setlength{\parskip}{0.6em}
\setstretch{1.08}
\setlist[itemize]{itemsep=0.4em}

% -----------------------------
% Farben (Corporate-Farbschema Physik KFR)
% -----------------------------
\definecolor{titleblue}{HTML}{183A6B}
\definecolor{accentblue}{HTML}{2E6F9E}
\definecolor{darkgray}{HTML}{4A4A4A}

% Links (hyperref): farblich ans Corporate-Farbschema angeglichen.
\hypersetup{
  colorlinks=true,
  urlcolor=accentblue,
  linkcolor=accentblue,
  pdftitle={Aufgabenblatt 3: Schaltungen},
  pdfauthor={Loris Keller}
}

% -----------------------------
% Boxen
% -----------------------------
\tcbset{before skip=10pt, after skip=10pt}
\newtcolorbox{titlebox}{
  enhanced,
  colback=titleblue,
  colframe=titleblue,
  boxrule=0pt,
  arc=3mm,
  left=3mm,right=3mm,top=3mm,bottom=3mm
}

\newlength{\aufgabenindent}
\settowidth{\aufgabenindent}{10.\hskip\labelsep}

\newtcolorbox{taskbox}[1][]{
  enhanced, breakable,
  colback=white, colframe=white, boxrule=0pt,
  left=0mm,right=0mm,top=0mm,bottom=0mm,
  before skip=10pt, after skip=10pt,
  before upper={%
    \def\tmpboxtitle{#1}%
    \ifx\tmpboxtitle\empty\else\noindent\textbf{#1}\par\smallskip\fi
  }
}

\newtcolorbox{answerbox}{
  breakable,
  colback=white,
  colframe=gray!55,
  boxrule=0.5pt,
  arc=1.5mm,
  left=2mm,right=2mm,top=2mm,bottom=2mm
}

% --- Loesungen ein-/ausblenden -----------------------------------------------
\noprintanswers
%\printanswers

\pointformat{}
\renewcommand{\solutiontitle}{\noindent\color{titleblue}\textbf{L\"osung:}\enspace}

\pagestyle{headandfoot}
\cfoot{\textcolor{darkgray}{Seite \thepage\ von \pageref{LastPage}}}

\newcommand{\Kopfzeile}[4]{%
  \begin{titlebox}
    {\color{white}\sffamily\bfseries\Large #4}
  \end{titlebox}
  \vspace{0.5em}
  \noindent\textbf{Klasse:}~#1 \hfill \textbf{Datum:}~#2 \hfill \textbf{Thema:}~#3
    \hfill \textbf{Lehrperson:}~Loris Keller
  \vspace{0.8em}
  \lhead{\textcolor{darkgray}{#3}}
  \rhead{\textcolor{darkgray}{#4}}
}

\newlength{\antwortraumlen}
\newcommand{\Antwortraum}[1]{%
  \pgfmathsetlength{\antwortraumlen}{#1*2.5cm}%
  \begin{answerbox}
    \rule{0pt}{\antwortraumlen}%
  \end{answerbox}%
}

% Werte fuer Kopfzeile zentral setzen
\newcommand{\Klasse}{1f}
\newcommand{\Datum}{23.09.2026}
\newcommand{\Thema}{Schaltungen}
\newcommand{\ArbeitsblattTitel}{Aufgabenblatt 3: Schaltungen}

\begin{document}

\Kopfzeile{\Klasse}{\Datum}{\Thema}{\ArbeitsblattTitel}

% =============================================================================
% Aufgabe 1: Zweite Lampe im Stromkreis (qualitativ)
% =============================================================================
\begin{questions}
\begin{taskbox}[Aufgabe 1: Zweite Lampe im Stromkreis]
\question[4] An eine Spannungsquelle ist zunächst eine einzelne Glühbirne
angeschlossen. Sie leuchtet in normaler Helligkeit.
\begin{center}
\begin{circuitikz}[scale=0.9]
  \draw (0,0) to[battery1,invert,l=$U_0$] (0,2.5)
              -- (3,2.5)
              to[lamp] (3,0)
              -- (0,0);
\end{circuitikz}
\end{center}
\begin{parts}
  \part[2] Was ändert sich, wenn eine zweite, baugleiche Glühbirne in
  \textbf{\emph{Serie}} zur bereits vorhandenen geschaltet wird?
  Begründen Sie ausführlich, auch bezüglich der Stromstärke, die die
  Spannungsquelle jetzt liefern muss.
  \begin{answerbox}
  \vspace{2.2cm}
  \end{answerbox}
  \begin{solution}
  Die beiden Glühbirnen liegen jetzt in Serie: Der Gesamtwiderstand der
  Schaltung ist die Summe der beiden (gleich grossen) Einzelwiderstände,
  er verdoppelt sich also. Da die Spannung der Quelle gleich bleibt,
  sinkt der Strom (Ohmsches Gesetz $I=U/R$). Beide Glühbirnen leuchten
  gleich hell, jedoch schwächer als die ursprüngliche einzelne
  Glühbirne vorher. Die Spannungsquelle muss dabei \textbf{\emph{weniger}}
  Strom liefern als zuvor.
  \end{solution}

  \part[2] Was ändert sich, wenn eine zweite, baugleiche Glühbirne
  stattdessen \textbf{\emph{parallel}} zur vorhandenen geschaltet wird?
  Begründen Sie wieder ausführlich, auch bezüglich der Stromstärke, die
  die Spannungsquelle liefern muss.
  \begin{answerbox}
  \vspace{2.2cm}
  \end{answerbox}
  \begin{solution}
  Bei einer Parallelschaltung liegt an jeder der beiden Glühbirnen
  weiterhin die volle Spannung der Quelle an, genau wie vorher an der
  einzelnen Glühbirne. Beide leuchten deshalb gleich hell wie die
  ursprüngliche einzelne Glühbirne, ihre Helligkeit ändert sich nicht.
  Der Gesamtstrom, den die Spannungsquelle liefern muss, verdoppelt sich
  aber, da sich der Strom nun auf zwei parallele Zweige mit je demselben
  Strom wie vorher aufteilt.
  \end{solution}
\end{parts}
\end{taskbox}
\end{questions}

% =============================================================================
% Aufgabe 2: Widerstandskombination
% =============================================================================
\needspace{14cm}
\begin{questions}
\setcounter{question}{1}
\begin{taskbox}[Aufgabe 2: Widerstandskombination]
\question[5] Gegeben ist die folgende Schaltung mit
$R_1=\SI{50}{\ohm}$, $R_2=\SI{100}{\ohm}$, $R_3=\SI{200}{\ohm}$ und
$U_0=\SI{220}{\volt}$:
\begin{center}
\includegraphics[width=6.5cm]{Aufgaben/Aufgabe_2/widerstandskombination-schaltplan.pdf}
\end{center}
\begin{parts}
  \part[2] Berechnen Sie den Gesamtwiderstand $R_{ges}$ der Schaltung.
  \begin{answerbox}
  \vspace{2.5cm}
  \end{answerbox}
  \begin{solution}
  $R_1$ und $R_2$ liegen in Serie:
  \[
    R_{12} = R_1+R_2 = \SI{50}{\ohm}+\SI{100}{\ohm} = \SI{150}{\ohm}.
  \]
  Diese Serie liegt parallel zu $R_3$:
  \[
    \frac{1}{R_{ges}} = \frac{1}{R_{12}}+\frac{1}{R_3}
    = \frac{1}{\SI{150}{\ohm}}+\frac{1}{\SI{200}{\ohm}}
    = \frac{4}{\SI{600}{\ohm}}+\frac{3}{\SI{600}{\ohm}}
    = \frac{7}{\SI{600}{\ohm}}
  \]
  \[
    \Longrightarrow\quad R_{ges} = \frac{\SI{600}{\ohm}}{7} \approx \SI{85.7}{\ohm}.
  \]
  \end{solution}

  \part[3] Berechnen Sie die Stromstärke $I_1$ durch $R_1$ und $R_2$
  sowie die Stromstärke $I_2$ durch $R_3$.
  \begin{answerbox}
  \vspace{2.5cm}
  \end{answerbox}
  \begin{solution}
  Beide Zweige liegen an derselben Spannung $U_0$ (Parallelschaltung),
  mit dem Ohmschen Gesetz gilt für jeden Zweig einzeln:
  \[
    I_1 = \frac{U_0}{R_{12}} = \frac{\SI{220}{\volt}}{\SI{150}{\ohm}}
    \approx \SI{1.47}{\ampere}.
  \]
  \[
    I_2 = \frac{U_0}{R_3} = \frac{\SI{220}{\volt}}{\SI{200}{\ohm}}
    = \SI{1.1}{\ampere}.
  \]
  \end{solution}
\end{parts}
\end{taskbox}
\end{questions}

% =============================================================================
% Aufgabe 3: Minimaler und maximaler Widerstand
% =============================================================================
\needspace{16cm}
\begin{questions}
\setcounter{question}{2}
\begin{taskbox}[Aufgabe 3: Minimaler und maximaler Widerstand]
\question[6] Gegeben ist die folgende Schaltung mit $U_0=\SI{12.0}{\volt}$,
$R_1=\SI{100}{\ohm}$ und $R_3=\SI{80.0}{\ohm}$. Die Grösse von $R_2$ ist
zunächst unbekannt.
\begin{center}
\includegraphics[width=6cm]{Aufgaben/Aufgabe_0/minimaler-maximaler-widerstand-schaltplan.pdf}
\end{center}
\begin{parts}
  \part[2] Begründen Sie ohne Rechnung, warum der Gesamtwiderstand
  dieser Schaltung grösser als $\SI{80}{\ohm}$, aber kleiner als
  $\SI{180}{\ohm}$ sein muss.
  \begin{answerbox}
  \vspace{2.5cm}
  \end{answerbox}
  \begin{solution}
  Die Schaltung ist eine Serieschaltung aus einer Parallelschaltung
  ($R_1\|R_2$) und dem Einzelwiderstand $R_3$. Da $R_3=\SI{80}{\ohm}$
  ist, muss der Gesamtwiderstand mindestens $\SI{80}{\ohm}$ betragen: Er
  kann durch die dazu in Serie liegende Parallelschaltung nur noch
  grösser werden. Der Widerstand einer Parallelschaltung ist aber immer
  kleiner als ihr kleinster Einzelwiderstand, hier also kleiner als
  $R_1=\SI{100}{\ohm}$. Der Gesamtwiderstand ist deshalb grösser als
  $\SI{80}{\ohm}$, aber kleiner als $\SI{80}{\ohm}+\SI{100}{\ohm}=\SI{180}{\ohm}$.
  \end{solution}

  \needspace{10cm}
  \part[2] Nun beträgt $R_2=\SI{60.0}{\ohm}$. Berechnen Sie die
  Stromstärke, die durch die Schaltung fliesst.
  \begin{answerbox}
  \vspace{2.5cm}
  \end{answerbox}
  \begin{solution}
  Zuerst $R_1$ parallel zu $R_2$:
  \[
    \frac{1}{R_{12}} = \frac{1}{R_1}+\frac{1}{R_2}
    = \frac{1}{\SI{100}{\ohm}}+\frac{1}{\SI{60}{\ohm}}
    = \frac{3}{\SI{300}{\ohm}}+\frac{5}{\SI{300}{\ohm}}
    = \frac{8}{\SI{300}{\ohm}}
  \]
  \[
    \Longrightarrow\quad R_{12} = \frac{\SI{300}{\ohm}}{8} = \SI{37.5}{\ohm}.
  \]
  Damit der Gesamtwiderstand (Serie mit $R_3$):
  \[
    R_{ges} = R_{12}+R_3 = \SI{37.5}{\ohm}+\SI{80}{\ohm} = \SI{117.5}{\ohm}.
  \]
  Mit dem Ohmschen Gesetz:
  \[
    I = \frac{U_0}{R_{ges}} = \frac{\SI{12.0}{\volt}}{\SI{117.5}{\ohm}}
    \approx \SI{0.1021}{\ampere}.
  \]
  \end{solution}

  \needspace{8cm}
  \part[2] Berechnen Sie die Spannungen, die an $R_1$, $R_2$ und $R_3$
  jeweils abfallen.
  \begin{answerbox}
  \vspace{2.5cm}
  \end{answerbox}
  \begin{solution}
  Für $R_3$ gilt mit dem Ohmschen Gesetz:
  \[
    U_3 = R_3\cdot I = \SI{80.0}{\ohm}\cdot\SI{0.1021}{\ampere}
    \approx \SI{8.17}{\volt}.
  \]
  An $R_1$ und $R_2$ fällt dieselbe Spannung ab, da sie parallel
  geschaltet sind. Diese ergibt sich aus der restlichen Spannung:
  \[
    U_1 = U_2 = U_0-U_3 = \SI{12.0}{\volt}-\SI{8.17}{\volt}
    \approx \SI{3.83}{\volt}.
  \]
  \end{solution}
\end{parts}
\end{taskbox}
\end{questions}

% =============================================================================
% Aufgabe 4: Widerstandsänderung
% =============================================================================
\needspace{16cm}
\begin{questions}
\setcounter{question}{3}
\begin{taskbox}[Aufgabe 4: Widerstandsänderung]
\question[5] Gegeben ist die folgende Schaltung mit $U=\SI{12}{\volt}$,
$R_1=\SI{10}{\ohm}$, $R_2=\SI{40}{\ohm}$, $R_3=\SI{20}{\ohm}$ und
$R_4=\SI{20}{\ohm}$:
\begin{center}
\includegraphics[width=6cm]{Aufgaben/Aufgabe_3/widerstandsaenderung-schaltplan.pdf}
\end{center}
\begin{parts}
  \part[3] Berechnen Sie den Gesamtwiderstand $R_{ges}$ und die
  Stromstärke an der Stelle $A$.
  \begin{answerbox}
  \vspace{3cm}
  \end{answerbox}
  \begin{solution}
  Die Schaltung ist eine Serieschaltung der beiden Parallelschaltungen
  $(R_1\|R_2)$ und $(R_3\|R_4)$. Zuerst der obere Teil:
  \[
    \frac{1}{R_{oben}} = \frac{1}{R_1}+\frac{1}{R_2}
    = \frac{1}{\SI{10}{\ohm}}+\frac{1}{\SI{40}{\ohm}}
    = \frac{4}{\SI{40}{\ohm}}+\frac{1}{\SI{40}{\ohm}}
    = \frac{5}{\SI{40}{\ohm}}
  \]
  \[
    \Longrightarrow\quad R_{oben} = \frac{\SI{40}{\ohm}}{5} = \SI{8}{\ohm}.
  \]
  Dann der untere Teil, auf die gleiche Weise:
  \[
    \frac{1}{R_{unten}} = \frac{1}{R_3}+\frac{1}{R_4}
    = \frac{1}{\SI{20}{\ohm}}+\frac{1}{\SI{20}{\ohm}}
    = \frac{2}{\SI{20}{\ohm}}
    \quad\Longrightarrow\quad R_{unten} = \SI{10}{\ohm}.
  \]
  Der Gesamtwiderstand ist die Summe der beiden (Serie):
  \[
    R_{ges} = R_{oben}+R_{unten} = \SI{8}{\ohm}+\SI{10}{\ohm} = \SI{18}{\ohm}.
  \]
  Mit dem Ohmschen Gesetz ergibt sich der Strom bei $A$:
  \[
    I = \frac{U}{R_{ges}} = \frac{\SI{12}{\volt}}{\SI{18}{\ohm}}
    \approx \SI{0.67}{\ampere}.
  \]
  \end{solution}

  \part[2] Nun wird $R_1$ durch einen $\SI{100}{\ohm}$-Widerstand
  ersetzt. Erläutern Sie \textbf{\emph{ohne Rechnung}}, ob die bei $B$
  gemessene Stromstärke danach kleiner, grösser oder gleich gross ist
  wie zuvor.
  \begin{answerbox}
  \vspace{2.5cm}
  \end{answerbox}
  \begin{solution}
  Durch den grösseren Widerstand an der Stelle $R_1$ wird der
  Gesamtwiderstand der oberen Parallelschaltung $R_{oben}$ grösser. Da
  diese in Serie mit der unteren Parallelschaltung liegt, wird dadurch
  auch der Gesamtwiderstand $R_{ges}$ grösser, und der Gesamtstrom
  sinkt. Da in einer Serieschaltung überall derselbe Strom fliesst,
  fliesst dieser kleinere Strom auch durch den unteren Teil der
  Schaltung, zu dem $B$ gehört: Die Stromstärke bei $B$ ist also
  \textbf{\emph{kleiner}} als zuvor.
  \end{solution}
\end{parts}
\end{taskbox}
\end{questions}

\end{document}
